\documentclass[12pt]{article}

% ── Packages ────────────────────────────────────────────────────────────────
\usepackage[margin=1in]{geometry}
\usepackage{setspace}
\doublespacing
\usepackage{amsmath,amssymb}
\usepackage{booktabs}
\usepackage{natbib}
\bibpunct{(}{)}{;}{a}{,}{,}
\usepackage{hyperref}
\hypersetup{colorlinks=true, linkcolor=black, citecolor=black, urlcolor=black}
\usepackage{graphicx}
\usepackage{caption}
\usepackage{threeparttable}
\usepackage{array}
\usepackage{microtype}

% ── Title ───────────────────────────────────────────────────────────────────
\title{\textbf{Appropriability and the Market Pricing of AI Model Capability}%
\thanks{We thank [ACKNOWLEDGMENTS]. All errors are our own.}}
\author{Zhuo Chen \and Ruishuang Lu \\[0.5em]
\small [Shandong University]}
\date{\small This version: June 2026}

% ────────────────────────────────────────────────────────────────────────────
\begin{document}
\maketitle
\thispagestyle{empty}

% ── Abstract ────────────────────────────────────────────────────────────────
\begin{abstract}
\noindent
Do financial markets price the capability of large language models (LLMs), or do
they merely react to AI hype? We study 60 major LLM release events from 2024 to
early 2026, constructing a firm-event panel of 86 AI-related listed companies. We
measure model capability using the Artificial Analysis Intelligence Index and
estimate cumulative abnormal returns over a 20-trading-day window. Closed-source
model releases drive a positive capability-pricing effect: a one-standard-deviation
increase in intelligence is associated with a 3.0 percentage point higher
CAR$[0,{+}20]$ ($\hat{\beta} = 0.0023$, $p = 0.008$ under wild cluster bootstrap).
This effect is absent for open-source releases, where the capability slope is
attenuated and reverses sign in point estimates. Media sentiment, measured by
FinBERT, is negatively associated with post-release returns---consistent with
narrative-driven pre-pricing that is distinct from technological capability.
Markets appear to price AI capability primarily when it can be appropriated.
\end{abstract}

\medskip
\noindent\textit{JEL codes:} G14, O30, L86\\
\noindent\textit{Keywords:} large language models, event study, appropriability, abnormal returns, media sentiment

\newpage

% ── 1. Introduction ─────────────────────────────────────────────────────────
\section{Introduction}
\label{sec:intro}

Large language model releases have become routine technology shocks---roughly 30
major releases per year in 2024--2025---attracting widespread investor attention.
Less clear is whether markets price the \emph{capability} embedded in these
releases or merely the \emph{narrative} surrounding them. This distinction matters
for understanding how financial markets process technological information and for
valuing firms whose competitive positions depend on appropriable AI performance.

Appropriability theory predicts that the market value of a technological advance
should depend on whether the innovating firm can capture the rents it generates
\citep{Teece1986}. Closed-source LLM deployments create proprietary
barriers---API lock-in, ecosystem advantages, and pricing power over model
access. Open-source releases, by contrast, diffuse capability broadly and may
commoditize inference services, eroding the competitive rents that would otherwise
be capitalized into stock prices. If appropriability determines market pricing,
the coefficient on model capability should be positive for closed-source releases
and weaker---or reversed---for open-source releases.

We test this prediction using a firm-event panel spanning 60 major LLM releases
from January 2024 through early 2026. We measure capability with the Artificial
Analysis (AA) Intelligence Index, a composite benchmark score that aggregates
performance across standardized reasoning and knowledge tasks. Our main outcome
is the market-model cumulative abnormal return over $[0,{+}20]$ trading days,
estimated separately for each firm-event pair and then regressed on the AA
Intelligence Index, conditioning on firm-level controls (size, book-to-market,
volatility, momentum) and year fixed effects. We cluster standard errors by event
and report results under both CR2 (Bell-McCaffrey small-cluster correction) and a
manually implemented Rademacher wild cluster bootstrap.

Three findings emerge. First, among closed-source releases, capability is
positively priced with strong statistical support: $\hat{\beta} = 0.0023$, CR2
$p = 0.007$, wild bootstrap $p = 0.008$, over 36 event clusters. A
one-standard-deviation increase in the Intelligence Index is associated with a
3.0 percentage-point higher 20-day abnormal return. Second, when we interact
capability with an open-source indicator, the slope for open-source releases is
substantially attenuated in point estimates (implied slope $= -0.0015$), though
the interaction term is only marginally significant under the most conservative
inference (CR2 $p = 0.048$, wild bootstrap $p = 0.086$); results should be read
as suggestive. Third, media sentiment measured by FinBERT is \emph{negatively}
associated with CAR$[0,{+}20]$ ($\hat{\beta} = -0.081$, CR2 $p = 0.003$), and
intelligence remains positively predictive after controlling for sentiment,
suggesting the capability effect reflects technological information rather than
narrative heat.

This paper connects to three bodies of work. Event studies on technology shocks
document positive abnormal returns around IT capability announcements
\citep[CITATION]{}), but they cannot separate capability from attention because
both co-move with the event. \citet{HallJaffeTrajtenberg2005} show that patent
citations---a proxy for technological appropriability---predict firm market value,
establishing that rent-capture moderates the value of innovation.
\citet{Tetlock2007} and \citet{EngelbergParsons2011} document that media sentiment
moves short-run prices; however, they do not test whether fundamental capability
survives sentiment controls or examine the direction of long-run returns. No study
we are aware of isolates \emph{which dimension} of AI model performance drives
post-release stock returns or whether the capability-to-value link depends on the
appropriability of the release. We provide this evidence, using direct capability
scores rather than proxies, and decompose LLM event returns into a
fundamental-capability component and a narrative component that move in opposite
directions.

The remainder of this paper proceeds as follows. Section~\ref{sec:data} describes
the data and variable construction. Section~\ref{sec:hyp} develops two testable
hypotheses. Section~\ref{sec:results} presents the main results.
Section~\ref{sec:robust} reports robustness checks. Section~\ref{sec:conclusion}
concludes.

% ── 2. Data ─────────────────────────────────────────────────────────────────
\section{Data and Variable Construction}
\label{sec:data}

\paragraph{Event sample.}
We compile 60 major LLM release events from January 2024 through March 2026,
drawing on model release announcements documented in AI benchmarking databases
and verified against official developer communications. Events are aggregated to
the model-family level when a developer releases multiple variants simultaneously;
a release event is included only if it corresponds to a meaningfully new
capability level rather than a minor update. The result is one observation per
release, with a confirmed public announcement date.

\paragraph{Firm-event panel.}
For each event we include up to 86 AI-related listed firms whose business
activities are plausibly connected to LLM development or deployment. Connection
types include direct model development (owner), equity investment in the developer
(investor), cloud infrastructure provision (cloud), upstream hardware supply,
downstream software integration, and direct competition. We identify these
relationships through public filings, investment records, and product
documentation, and we code them as binary pre-event indicators. We control for
firm-model exposure using these pre-event relationship indicators; full variable
definitions are in Appendix~A. The full panel has 5,161 firm-event observations;
the regression sample with non-missing Intelligence Index and controls contains
3,780 observations across 47 event clusters.

\paragraph{Cumulative abnormal returns.}
We estimate daily abnormal returns using a market model with parameters estimated
over a $[-200, -21]$ pre-event window \citep{MacKinlay1997}. We then accumulate
abnormal returns over post-announcement windows, with the primary window
$[0,{+}20]$ (twenty trading days). The event date is the verified public
announcement date.

\paragraph{Model capability.}
The AA Intelligence Index is a composite score produced by Artificial Analysis
that aggregates performance across mathematics, reasoning, and general knowledge
benchmarks. It has a mean of 26.6 and a standard deviation of 13.1 in our
regression sample. We center the index at its sample mean ($\textit{intel\_c}$)
so that main effects of binary moderators are interpretable at average capability
levels.

\paragraph{Open-source indicator.}
We code $\textit{is\_open\_weight} = 1$ if the model is released under an
open-weight or fully open-source license permitting free redistribution and
inference, and 0 otherwise. Among the 60 events, 11 are open-source (881
firm-event observations) and 36 have non-missing intelligence data under a
closed-source designation (2,902 observations).

\paragraph{Media sentiment.}
For each event we compute the mean FinBERT-based sentiment score of news articles
published in a $\pm 5$-day window around the announcement date. The sentiment
variable has a mean of 0.18 and a standard deviation of 0.16, indicating that AI
model releases attract generally positive coverage. We center sentiment at its
sample mean ($\textit{sent\_c5}$).

\paragraph{Controls.}
Firm-level controls are the log of total assets, the book-to-market ratio,
realized return volatility over the prior quarter, and momentum (prior
12-minus-1-month return). Year fixed effects absorb aggregate time trends. All
controls are measured as of the quarter preceding the event.

% ── 3. Hypotheses ───────────────────────────────────────────────────────────
\section{Hypothesis Development}
\label{sec:hyp}

Appropriability theory \citep{Teece1986} holds that a technological advance
creates stock market value only to the extent that the innovating firm---or firms
positioned to benefit---can appropriate the rents it generates. Applied to LLM
releases, this logic yields two testable predictions.

\medskip
\noindent\textbf{H1} (Closed-source capability pricing). \textit{Model capability
is positively associated with post-release cumulative abnormal returns when the
model is closed-source.} A higher-capability closed-source model raises the
prospect of expanded API revenue, stronger customer lock-in, and a wider
competitive moat.

\medskip
\noindent\textbf{H2} (Open-source attenuation). \textit{The capability-pricing
slope is weaker for open-source releases than for closed-source releases.}
Open-weight models allow unrestricted use and redistribution, eroding proprietary
barriers. Note that the prediction concerns \emph{relative slopes}, not the
unconditional average return.

\medskip
We additionally examine whether media sentiment explains the capability result.
If intelligence merely reflects AI hype, controlling for sentiment should absorb
the intelligence coefficient. The alternative---that intelligence reflects
technological information beyond sentiment---implies that both retain independent
predictive power.

% ── 4. Results ──────────────────────────────────────────────────────────────
\section{Results}
\label{sec:results}

\subsection{Baseline capability pricing}

Table~\ref{tab:baseline} presents the baseline results. Among closed-source
releases, capability is strongly and precisely priced: $\hat{\beta} = 0.0023$
(SE $= 0.0006$), CR2 $p = 0.007$, wild bootstrap $p = 0.008$, over 36 event
clusters. \textbf{A one-standard-deviation increase in the AA Intelligence Index
(13.1 points) is associated with a 3.0-percentage-point higher 20-day abnormal
return for firms exposed to a closed-source release.} The result survives year
fixed effects and the full set of firm-level controls, consistent with H1.

The full-sample estimate is positive but weaker: $\hat{\beta} = 0.0015$
(SE $= 0.0007$), CR0 $p = 0.027$, CR2 $p = 0.054$, wild bootstrap $p = 0.053$.
The full-sample coefficient falls short of conventional thresholds once
small-cluster corrections are applied---a natural consequence of pooling events
where the capability channel is active (closed-source) with those where it is
attenuated (open-source).

The contrast with open-source releases supports H1: the appropriable capability
channel operates specifically in settings where proprietary access is maintained.

\subsection{Open-source attenuation}

Table~\ref{tab:interaction} examines whether open-source status moderates
capability pricing by estimating:
\begin{equation}
\label{eq:interaction}
\text{CAR}[0,{+}20]_{it} = \beta_1\,\textit{intel\_c}_{e}
  + \beta_2\,\textit{open}_{e}
  + \beta_3\,(\textit{intel\_c}_{e} \times \textit{open}_{e})
  + \mathbf{X}_{it}'\boldsymbol{\gamma}
  + \alpha_t + \varepsilon_{it}
\end{equation}
where $e$ indexes event, $i$ firm, $t$ year, and $\mathbf{X}_{it}$ contains
firm-level controls.

The baseline (closed-source) slope $\hat{\beta}_1 = +0.0023$ (CR2 $p = 0.005$).
The interaction coefficient $\hat{\beta}_3 = -0.0037$ (SE $= 0.0012$), implying
a substantially attenuated open-source slope of $-0.0015$. The interaction is
significant under CR0 ($p = 0.003$) and marginally so under CR2 ($p = 0.048$);
under the wild cluster bootstrap, however, $p = 0.086$, short of the conventional
5\% threshold.

We interpret this evidence cautiously. The point estimates are consistent with
H2---open-source status attenuates and in point estimates reverses the
capability-pricing slope---but only 11 open-source events identify the interaction,
and the wild bootstrap signals fragile inference at this cluster count. We treat
this as suggestive evidence consistent with the appropriability mechanism, rather
than a robust finding.

The open-source main effect at average intelligence ($\hat{\beta}_2 = -0.013$) is
not statistically significant, consistent with the interaction story: open-source
releases do not unconditionally yield lower returns; they yield lower
\emph{marginal returns to additional capability}.

\subsection{Media sentiment and the hype alternative}

Table~\ref{tab:sentiment} addresses whether the capability effect is confounded
with media hype. Column~(1) shows that sentiment alone strongly and
\emph{negatively} predicts CAR$[0,{+}20]$: $\hat{\beta} = -0.081$,
CR2 $p = 0.003$. More positive media coverage in the event window is associated
with \emph{lower} subsequent 20-day returns. One interpretation is that events
generating strong positive narratives have already experienced anticipatory price
appreciation, so the $[0,{+}20]$ abnormal return is depressed by a prior run-up
that our pre-event estimation window partially misses. We cannot identify the
precise mechanism with the current data.

Columns~(2)--(3) show the joint model. Intelligence and sentiment both retain
their sign and significance under CR0; under CR2, the intelligence coefficient is
marginal ($p = 0.087$) and sentiment is marginal ($p = 0.057$). Crucially,
controlling for sentiment does not eliminate the intelligence coefficient---the
two variables are orthogonal at the event level (event-level regression of
sentiment on intelligence: $p = 0.317$), confirming they capture distinct
information. The interaction between intelligence and sentiment is not significant
($\hat{\beta} = +0.002$, $p = 0.555$), so sentiment neither amplifies nor dampens
the capability channel.

% ── 5. Robustness ───────────────────────────────────────────────────────────
\section{Robustness}
\label{sec:robust}

\paragraph{Alternative windows.}
Appendix Table~A1 repeats the closed-source regression across windows from
CAR$[0,{+}1]$ to CAR$[0,{+}20]$. The coefficient is largest and most precisely
estimated at the long window; the immediate $[0,{+}1]$ response is positive but
not significant, suggesting capability information takes time to be fully
processed. This pattern is consistent with fundamental revaluation rather than
noise trading around announcements.

\paragraph{Fama-French three-factor model.}
Replacing the market model with the Fama-French three-factor model leaves the
main coefficients unchanged (Appendix Table~A1).

\paragraph{Industry and Mag7 checks.}
Internet infrastructure, software, and semiconductor firms show the strongest
positive response (CR2 $p \leq 0.015$ for the internet infrastructure subgroup).
Results are not driven by the Magnificent Seven---the interaction between
intelligence and a Mag7 indicator is not significant ($p = 0.633$), and the
capability slope is nearly identical for Mag7 and non-Mag7 firms (Appendix
Table~A2).

\paragraph{Owner short-window effects.}
Firms that are direct owners of the releasing model show a negative short-window
abnormal return (CAR$[0,{+}1]$: $\hat{\beta} = -0.0018$ for the
$\textit{intel\_c} \times \textit{owner}$ interaction, wild $p = 0.034$). This
pattern---in which the releasing firm's equity does not benefit immediately---may
reflect profit-taking or uncertainty about commercial adoption timelines. We treat
this as suggestive and report full details in Appendix Table~A3.

% ── 6. Conclusion ────────────────────────────────────────────────────────────
\section{Conclusion}
\label{sec:conclusion}

Financial markets price AI model capability selectively. Among closed-source
releases, a one-standard-deviation increase in the Artificial Analysis
Intelligence Index is associated with a 3.0-percentage-point higher 20-day
cumulative abnormal return---a result that is robust to conservative small-cluster
inference. Among open-source releases, the capability slope is attenuated and
reverses sign in point estimates, though this interaction rests on only 11
open-source event clusters and should be treated as exploratory.

Media sentiment is negatively associated with post-release returns and operates
independently of the capability channel, consistent with a mechanism in which
positive narratives bring forward price gains that subsequently dissipate.
Intelligence and sentiment are jointly significant in the full model, providing a
rough decomposition of LLM event returns into a fundamental-capability component
and a narrative component that move in opposite directions.

Financial markets appear to price the \emph{economic ownership} of AI performance,
not AI performance itself. As open-source LLMs mature and inference costs fall,
the appropriability channel we identify may weaken, compressing the abnormal
returns to AI capability improvements among affected firms.

% ── References ───────────────────────────────────────────────────────────────
\newpage
\bibliographystyle{aer}
\begin{thebibliography}{9}

\bibitem[Engelberg and Parsons(2011)]{EngelbergParsons2011}
Engelberg, Joseph~E., and Christopher~A. Parsons. 2011.
``The Causal Impact of Media in Financial Markets.''
\textit{Journal of Finance} 66(1): 67--97.

\bibitem[Hall, Jaffe, and Trajtenberg(2005)]{HallJaffeTrajtenberg2005}
Hall, Bronwyn~H., Adam Jaffe, and Manuel Trajtenberg. 2005.
``Market Value and Patent Citations.''
\textit{RAND Journal of Economics} 36(1): 16--38.

\bibitem[MacKinlay(1997)]{MacKinlay1997}
MacKinlay, A.~Craig. 1997.
``Event Studies in Economics and Finance.''
\textit{Journal of Economic Literature} 35(1): 13--39.

\bibitem[Teece(1986)]{Teece1986}
Teece, David~J. 1986.
``Profiting from Technological Innovation: Implications for Integration,
Collaboration, Licensing and Public Policy.''
\textit{Research Policy} 15(6): 285--305.

\bibitem[Tetlock(2007)]{Tetlock2007}
Tetlock, Paul~C. 2007.
``Giving Content to Investor Sentiment: The Role of Media in the Stock Market.''
\textit{Journal of Finance} 62(3): 1139--1168.

% [CITATION: technology shocks and asset pricing — add appropriate reference]
% [CITATION: AI/LLM event studies — add concurrent working papers when available]
% [CITATION: biotech parallel for owner effect — add appropriate reference]

\end{thebibliography}

% ── Tables ────────────────────────────────────────────────────────────────────
\newpage
\section*{Tables}

% ─ Table 1 ─
\begin{table}[htbp]
\centering
\caption{Baseline Effect of AI Model Capability on Cumulative Abnormal Returns}
\label{tab:baseline}
\begin{threeparttable}
\begin{tabular}{lcccc}
\toprule
 & (1) & (2) & (3) & (4) \\
 & Full sample & Closed-source & Full sample & Closed-source \\
\midrule
AA Intelligence Index (centered) & 0.00152* & 0.00232*** & 0.00152 & 0.00232** \\
 & (0.00067) & (0.00062) & & \\[4pt]
CR0 $p$-value & 0.027 & 0.001 & & \\
CR2 $p$-value & & & 0.054 & 0.007 \\
Wild bootstrap $p$-value & & & 0.053 & 0.008 \\[6pt]
\midrule
Year FE & Yes & Yes & Yes & Yes \\
Firm controls & Yes & Yes & Yes & Yes \\
Observations & 3,780 & 2,899 & 3,780 & 2,899 \\
Event clusters & 47 & 36 & 47 & 36 \\
\bottomrule
\end{tabular}
\begin{tablenotes}[flushleft]\footnotesize
\item Dependent variable is the market-model CAR$[0,{+}20]$.
\textit{intel\_c} is the AA Intelligence Index centered at its sample mean (26.6).
Columns (1)--(2) report CR0 clustered standard errors in parentheses;
Columns (3)--(4) add CR2 and wild bootstrap $p$-values.
Firm controls are log total assets, book-to-market, prior-quarter volatility, and
12-minus-1-month momentum.
Wild bootstrap uses $B = 4{,}999$ Rademacher draws with the impose-null approach.
$+\,p<0.10$; $^{*}\,p<0.05$; $^{**}\,p<0.01$; $^{***}\,p<0.001$.
\end{tablenotes}
\end{threeparttable}
\end{table}

% ─ Table 2 ─
\begin{table}[htbp]
\centering
\caption{The Role of Appropriability: Open-Source Interaction}
\label{tab:interaction}
\begin{threeparttable}
\begin{tabular}{lcc}
\toprule
 & (1) CR0 & (2) CR2 \\
\midrule
\textit{intel\_c} (closed-source slope) & 0.00226*** & 0.00226** \\
 & (0.00060) & \\[2pt]
CR0 $p$ & 0.001 & \\
CR2 $p$ & & 0.005 \\[6pt]
\textit{is\_open\_weight} (at mean intelligence) & $-$0.013 & $-$0.013 \\
 & (0.012) & \\[6pt]
\textit{intel\_c} $\times$ \textit{is\_open\_weight} & $-$0.00373** & $-$0.00373* \\
 & (0.00120) & \\[2pt]
CR0 $p$ & 0.003 & \\
CR2 $p$ & & 0.048 \\
Wild bootstrap $p$ & 0.086 & 0.086 \\[6pt]
\textit{Implied open-source slope} ($\hat{\beta}_1 + \hat{\beta}_3$) & $-$0.00148 & $-$0.00148 \\[4pt]
\midrule
Year FE & Yes & Yes \\
Firm controls & Yes & Yes \\
Observations & 3,780 & 3,780 \\
Event clusters & 47 & 47 \\
\bottomrule
\end{tabular}
\begin{tablenotes}[flushleft]\footnotesize
\item Dependent variable is CAR$[0,{+}20]$. Estimation equation is
(\ref{eq:interaction}). CR0 standard errors in parentheses for Column (1).
The interaction term is significant under CR0 and marginally so under CR2;
the wild bootstrap $p$-value exceeds 0.05 and this result should be interpreted
with caution. Other notes as in Table~\ref{tab:baseline}.
\end{tablenotes}
\end{threeparttable}
\end{table}

% ─ Table 3 ─
\begin{table}[htbp]
\centering
\caption{Media Sentiment, Model Capability, and CAR$[0,{+}20]$}
\label{tab:sentiment}
\begin{threeparttable}
\begin{tabular}{lcccc}
\toprule
 & (1) & (2) & (3) & (4) \\
 & Sent. only & Intel. only & Joint & Joint (CR2) \\
\midrule
\textit{sent\_c5} (centered FinBERT) & $-$0.081*** & & $-$0.061* & $-$0.061+ \\
 & (0.020) & & (0.026) & \\[2pt]
CR2 $p$ (sent) & 0.003 & & & 0.057 \\[6pt]
\textit{intel\_c} & & 0.00152* & 0.00158* & 0.00158+ \\
 & & (0.00067) & (0.00072) & \\[2pt]
CR2 $p$ (intel) & & 0.054 & & 0.087 \\[6pt]
\midrule
Year FE & Yes & Yes & Yes & Yes \\
Firm controls & Yes & Yes & Yes & Yes \\
Observations & 3,068 & 3,780 & 3,068 & 3,068 \\
Event clusters & 47 & 47 & 47 & 47 \\
\bottomrule
\end{tabular}
\begin{tablenotes}[flushleft]\footnotesize
\item \textit{sent\_c5} is the FinBERT mean sentiment over a $\pm 5$-day event
window, centered at its sample mean (0.18). Sentiment is available for
approximately 3,068 observations. Column (4) reports CR2 $p$-values for both
predictors; CR0 standard errors are in parentheses for Columns (1)--(3).
The event-level correlation between \textit{intel\_c} and \textit{sent\_c5}
is not statistically significant ($p = 0.317$), confirming orthogonality.
Other notes as in Table~\ref{tab:baseline}.
\end{tablenotes}
\end{threeparttable}
\end{table}

\end{document}
